\scanpage{13}
- \mspage{10} -

(3.1), $c_4=0$ si et seulement si $a_2^2=0$. Supposons qu'on ait même $a_2=0$
\[
(5.1)\qquad y^2=x^3+a_4x+a_6\qquad\pi=a_4^{-1}dy=dx/2y
\]
\[
c_4=c_6=0\qquad\Delta=-a_4^3.
\]
Les autres systèmes $(\pi,x,y)$ donnant une équation de cette forme s'obtiennent en faisant $s=t=0$ dans (1.2) et
\[
\begin{aligned}
u^4a'_4&=a_4\\
(5.2)\qquad&\\
u^6a'_6&=a_6+ra_4+r^3.
\end{aligned}
\]
b) $p=2$

En caractéristique 2, on a $c_4=a_1^4$. Supposons qu'on ait $a_1=0$ ; $\pi$ étant donnée, on peut alors choisir $x$ et $y$ tels que (1.3) soit
\[
(5.3)\qquad y^2+a_3y=x^3+a_4x+a_6\qquad\pi=\frac{dx}{a_3}=dy/{}_{x^2+a_4}:
\]
\[
b_2=b_4=0\qquad b_6=a_3^2\qquad b_8=a_4^2\qquad\Delta=a_3^4\qquad c_4=c_6=0.
\]
Les autres systèmes $(\pi,x,y)$ donnant une équation de cette forme s'obtiennent en faisant $r=s^2$ dans (1.2) et
\[
\begin{aligned}
u^3\unsure{a'_3|e'_3}&=a_3\\
(5.4)\qquad u^4a'_4&=a_4+sa_3+s^4\\
u^6a'_6&=a_6+s^2a_4+ta_3+s^6+t^2.
\end{aligned}
\]
\margin{12}
